% =====================================================================
%  SECCION 6 -- PRODUCTO MINIMO VIABLE
%  SECCION 7 -- CONCLUSIONES, LIMITACIONES Y TRABAJO PENDIENTE
% =====================================================================

\newpage
\section{Producto Mínimo Viable (MVP)}

\subsection{Alcance del prototipo}

El MVP materializa el subconjunto de requisitos funcionales de prioridad \emph{Must} suficiente para
validar con personas usuarias reales la hipótesis de valor del sistema: que el registro digital de
la operación de campo y la centralización de alertas reducen la latencia entre la ocurrencia de un
problema y su atención.

El prototipo es una aplicación de interfaz construida con React, Vite y Tailwind CSS, con nueve
pantallas operativas, control de acceso por tres perfiles y persistencia local en el navegador. El
código fuente reside en un repositorio Git separado, enlazado desde \id{05\_MVP/README.md}.

\subsubsection{Cobertura verificada}

La siguiente tabla se construyó por inspección del código fuente, no por declaración del equipo de
desarrollo. El estado \emph{simulado} identifica funcionalidades cuya interfaz está completa pero
cuyo resultado no proviene de un procesamiento real.

\footnotesize
\begin{longtable}{|C{1.4cm}|L{5.6cm}|C{1.4cm}|C{2.6cm}|L{3.4cm}|}
\hline
\rowcolor{grisclaro}\textbf{ID-RF} & \textbf{Funcionalidad} & \textbf{WSJF} & \textbf{Estado} & \textbf{Pantalla} \\
\hline
\endhead
\id{RF-01} & Autenticación y control de acceso por rol & 4,20 & Implementado & Login \\ \hline
\id{RF-02} & Gestión de plantaciones y lotes & 3,67 & Implementado & Detalle de lote \\ \hline
\id{RF-03} & Gestión de personal y equipos & 3,67 & Implementado & Personal \\ \hline
\id{RF-04} & Registro de labores agrícolas & 4,20 & Implementado & Labores \\ \hline
\id{RF-05} & Registro de monitoreo fitosanitario & 4,80 & Implementado & Análisis \\ \hline
\id{RF-12} & Generación de alertas tempranas & 4,20 & Implementado & Alertas \\ \hline
\id{RF-19} & Generación de reportes & 3,20 & Implementado & Reportes \\ \hline
\id{RF-20} & Historial de actividades y eventos & 2,67 & Implementado & Detalle de lote \\ \hline
\id{RF-30} & Reporte de incidencia desde campo & 4,80 & Implementado & Análisis \\ \hline
\id{RF-13} & Registro del proceso de polinización & 4,20 & Parcial & Labores \\ \hline
\id{RF-15} & Visualización de recorridos del personal & 2,60 & Parcial & Mapa GPS \\ \hline
\id{RF-28} & Registro de avance por unidad de labor & 3,60 & Parcial & Labores \\ \hline
\id{RF-07} & Detección de plagas por imagen & 3,62 & \textbf{Simulado} & Análisis \\ \hline
\id{RF-14} & Conteo georreferenciado de flores & 3,25 & \textbf{Simulado} & Mapa GPS \\ \hline
\id{RF-08} & Diagnóstico nutricional por imagen & 2,25 & No implementado & --- \\ \hline
\id{RF-10} & Gestión de variedades y umbrales & 7,00 & No implementado & --- \\ \hline
\id{RF-18} & Estimación de producción y rendimiento & 3,60 & No implementado & --- \\ \hline
\id{RF-21} & Clasificación de madurez del racimo & 4,25 & No implementado & --- \\ \hline
\id{RF-22} & Alerta preventiva de fruta verde & 4,25 & No implementado & --- \\ \hline
\id{RF-26} & Planificación semanal con presupuesto & 3,60 & No implementado & --- \\ \hline
\id{RF-35} & Registro delegado del avance & 7,00 & No implementado & --- \\ \hline
\caption{Cobertura verificada del MVP sobre los requisitos funcionales \emph{Must}}
\label{tab:mvp}
\end{longtable}

\noindent
\textbf{Cobertura alcanzada.} De los 19 requisitos funcionales \emph{Must} vigentes al momento de
construcción del prototipo, 9 se encuentran implementados de forma completa, 3 de forma parcial y 2
en estado simulado. La cobertura efectiva es de $9/19 = 47{,}4\%$, o de $10{,}5/19 = 55{,}3\%$
contabilizando las implementaciones parciales con peso medio.

\noindent\fbox{\parbox{0.97\textwidth}{%
\small\textbf{Declaración de incumplimiento del mínimo.} La §3.6 de la guía exige una cobertura de
al menos el 60\% de los requisitos \emph{Debe tener}. El prototipo entregado no alcanza ese umbral.
Se declara de forma expresa el valor real (47{,}4\% completo, 55{,}3\% ponderado) en lugar de
computar como implementadas las funcionalidades simuladas, lo que habría permitido reportar un
73{,}7\% sin sustento verificable. Los requisitos ausentes y su plan de incorporación se detallan en
la §7.5. Esta limitación se registra como \id{L-10}.}}

\subsubsection{Funcionalidades simuladas}

Dos funcionalidades presentan interfaz completa con resultado no real, y así se declaran tanto en
este documento como en el archivo \id{README.md} del repositorio del prototipo:

\begin{itemize}[nosep]
  \item \textbf{Análisis de imagen (\id{RF-07}).} La pantalla captura la imagen, muestra el
        indicador de calidad y presenta un resultado de diagnóstico, pero dicho resultado se genera
        mediante una función pseudoaleatoria y no mediante inferencia sobre la imagen. La
        arquitectura de integración ---orquestador, interfaz del servicio y presentación de la
        explicación--- sí está construida.
  \item \textbf{Conteo georreferenciado (\id{RF-14}).} El mapa presenta recorridos y totales de
        polinización a partir de datos de ejemplo fijos; no se consulta el receptor GPS del
        dispositivo.
\end{itemize}

\subsection{Decisión de alcance: por qué la IA no se implementó}

La ausencia de inferencia real es deliberada y se argumenta.

Un clasificador de madurez o de plagas con utilidad real exige un conjunto de datos etiquetado del
propio cultivo. La literatura del dominio reporta conjuntos de entre 850 y varios miles de imágenes
para alcanzar exactitudes superiores al 90\% \cite{suharjito2021,omar2024,rosbi2024}. El equipo no
dispone de ese material: las fotografías recogidas en el trabajo de campo documentan el entorno y
las sesiones de elicitación, y no constituyen un conjunto etiquetado de entrenamiento.

Integrar un modelo preentrenado de propósito general habría producido una demostración visualmente
convincente pero sin validez: las predicciones no serían atribuibles al dominio y los valores
objetivo de \id{RNF-01} y \id{RNF-02} no podrían verificarse. Se prefiere entregar la arquitectura
de integración con la simulación \emph{declarada como tal} antes que presentar un diagnóstico
generado aleatoriamente como si fuera un resultado de inferencia.

\subsection{Desviaciones detectadas y plan de corrección}

La revisión del prototipo frente a la especificación identificó tres desviaciones que se declaran y
cuya corrección se programa antes de la Entrega 4.

\begin{longtable}{|C{1.4cm}|L{6.6cm}|L{7.0cm}|}
\hline
\rowcolor{grisclaro}\textbf{ID} & \textbf{Desviación} & \textbf{Corrección programada} \\
\hline
\endhead
\id{D-01} & El catálogo de datos del prototipo no corresponde al dominio real: emplea lotes con nomenclatura alfanumérica, cinco tipos de labor genéricos y nombres de personas ficticios, en lugar de la estructura de seis lotes, las labores documentadas en \id{EV-11} y \id{EV-13}, y los códigos de participante. & Sustituir los datos de ejemplo por la estructura real anonimizada: lotes 1 a 6, catálogo de labores del documento de liquidación (\id{RF-36}) y códigos \id{ENTR-XX} en lugar de nombres propios. \\ \hline
\id{D-02} & Las credenciales se almacenan en texto plano en el navegador, en contradicción con \id{RNF-13}. & Trasladar la autenticación a un servicio con derivación de clave; mientras tanto, se declara en el \id{README.md} del prototipo que las cuentas son de demostración y no deben emplearse con datos reales. \\ \hline
\id{D-03} & El código fuente se distribuye como archivo comprimido dentro del repositorio, lo que impide la revisión por archivo y el seguimiento de cambios. & Publicar el árbol de fuentes sin comprimir y declarar en el \id{package.json} el nombre del proyecto y su procedencia como exportación de herramienta de prototipado. \\ \hline
\caption{Desviaciones detectadas entre el prototipo y la especificación}
\end{longtable}

\subsection{Repositorio, despliegue y demostración}

El código fuente reside en un repositorio Git separado, enlazado desde \id{05\_MVP/README.md}
mediante submódulo, de modo que quede trazado el \emph{commit} exacto evaluado. El repositorio
contiene el código organizado por módulos, la tabla de cobertura de la Tabla~\ref{tab:mvp} y las
instrucciones de ejecución. El despliegue local requiere Node.js 18 o superior:

\begin{quote}
\id{npm install}\\
\id{npm run dev}
\end{quote}

\noindent
La persistencia se resuelve en el almacenamiento local del navegador. No existe todavía servicio de
respaldo ni base de datos, por lo que los requisitos no funcionales de fiabilidad
(\id{RNF-10}, \id{RNF-11}) y de seguridad (\id{RNF-13}, \id{RNF-14}) no son verificables sobre el
prototipo en su estado actual.

El recorrido funcional se documenta en \id{05\_MVP/video\_demo.mp4}, con duración no superior a tres
minutos, y cubre la secuencia: inicio de sesión por rol, consulta del plan semanal, registro de
labor con avance, registro delegado, reporte de incidencia con fotografía y coordenadas, y consulta
del consolidado.

% =====================================================================
\newpage
\section{Conclusiones, limitaciones y trabajo pendiente}

\subsection{Logros de la entrega}

La Entrega 3 (2A) consolida un documento de especificación completo sustentado en trabajo de campo
real. Los resultados verificables son los siguientes.

\textbf{Especificación.} Se pasó de 20 a 39 requisitos funcionales, de 9 a 18 requisitos no
funcionales cuantificados sobre las nueve características de ISO/IEC 25010:2023, de 7 a 10
restricciones de diseño, y se incorporaron 8 requisitos legales trazados a la LOPDP. Los 19
requisitos \emph{Must} cuentan con historia de usuario en formato Connextra, criterios INVEST
verificados y escenario de aceptación en Gherkin.

\textbf{Modelado.} Se produjeron trece diagramas: contexto, dos modelos organizacionales i*
(dependencia y razón estratégica), casos de uso general, clases refinado con atributos y
operaciones, tres de secuencia, actividad, dos de estados, componentes y despliegue. Se especificaron
diez casos de uso de forma textual, con al menos dos flujos alternativos y una excepción cada uno.

\textbf{Trazabilidad.} La matriz enlaza nueve niveles en 48 filas con trece columnas. Se corrigió la
deficiencia señalada en la Entrega 2, en la que dos requisitos funcionales carecían de caso de uso
asociado; la cobertura requisito-caso de uso es ahora del 100\%.

\textbf{Campo.} Se ejecutó una segunda ronda con cinco nuevos participantes, incorporando por
primera vez una organización externa ---la planta extractora--- como fuente de requisitos. El
cuestionario se amplió de 4 a 62 respondientes. Se obtuvieron dos tipos documentales de la
organización: el plan semanal de labores (\id{EV-11}) y la hoja de liquidación de presupuesto contra
ejecución (\id{EV-13}), esta última fuente de cuatro requisitos funcionales nuevos.

\subsection{Hallazgos que modificaron decisiones previas}

Tres resultados de la segunda ronda invalidaron supuestos que la Entrega 2 daba por establecidos.
Se registran de forma explícita porque ilustran el valor de la triangulación.

\textbf{Primero: el umbral de madurez no es constante.} La primera fuente de la extractora
(\id{EV-08}) indicó un criterio de cinco frutos desprendidos. La segunda (\id{EV-04}) precisó que el
umbral depende de la variedad: tres frutos en \emph{guineensis} y cinco en material híbrido. De
haberse especificado \id{RF-21} con una sola fuente, el requisito habría quedado incorrectamente
fijado en un valor constante. Este hallazgo motivó la restricción \id{RD-08} y el requisito
\id{RNF-18} de parametrización.

\textbf{Segundo: no todo el personal dispone de teléfono inteligente.} La Entrega 2 asumía lo
contrario a partir de cuatro respondientes. Con $n=62$, el 11{,}3\% declara no usarlo. La suposición
se retiró y se incorporaron \id{RF-35} y \id{RD-10} para habilitar el registro delegado. Sin esa
corrección, aproximadamente una de cada nueve personas destinatarias habría quedado excluida del
cómputo de avance y, por tanto, de su remuneración por rendimiento.

\textbf{Tercero: el rastreo por GPS no goza de aceptación unánime.} La aplicación piloto sugería
una aceptación del 75\%, que se leyó como conformidad. Con $n=62$, el 74{,}2\% no manifiesta
preocupación, pero el 25{,}8\% restante sí expresa reservas. En consecuencia, \id{RL-03} exige
consentimiento específico y revocable sin consecuencia laboral, y \id{CU-03} incorpora una rama que
permite ejecutar la labor sin rastreo.

\textbf{Cuarto: la tarifa de cada labor ya está formalizada en la organización.} La
especificación asumía que el valor unitario del trabajo por avance podía no estar configurado, y así
constaba como excepción del caso de uso \id{CU-08}. La hoja de liquidación semanal (\id{EV-13})
demuestra que la organización mantiene un catálogo de labores con código, unidad de medida y tarifa
vigente, y que liquida semanalmente el presupuesto contra la ejecución con aprobación de la
administración. El sistema no debe crear ese modelo: debe reproducir el que ya existe. De ahí los
requisitos \id{RF-36} a \id{RF-39}.

Este hallazgo ilustra una limitación de la entrevista como técnica única: ninguno de los ocho
participantes describió el catálogo de tarifas, pese a que todos operan bajo él. Fue necesario el
análisis documental para hacerlo explícito, lo que confirma el valor de la triangulación entre
técnicas y no solo entre fuentes.

\subsection{Limitaciones declaradas}

Se declaran de forma expresa las limitaciones de esta entrega. Ninguna se presenta como resuelta.

\begin{longtable}{|C{1.2cm}|L{6.4cm}|L{7.4cm}|}
\hline
\rowcolor{grisclaro}\textbf{ID} & \textbf{Limitación} & \textbf{Mitigación aplicada y estado} \\
\hline
\endhead
\id{L-01} & El perfil dominante del cuestionario ($n=60$) agrupa cuatro sub-perfiles de labor cuyos estratos individuales no alcanzan los 30 respondientes. & Se declara el agrupamiento en la §2.8 y el análisis por estrato se reporta como exploratorio, no inferencial. \\ \hline
\id{L-02} & La pregunta sobre comodidad con aplicaciones registró un 40{,}3\% de no respuesta. & El valor se presenta como indicio descriptivo y no se emplea para inferencia. Se revisará el instrumento antes de la siguiente aplicación. \\ \hline
\id{L-03} & La segunda ronda incorporó cinco entrevistas, por debajo del objetivo de ocho de la §4.2 de la guía. & Se declara el número real. La saturación temática se reporta en el Apéndice B con la curva correspondiente. \\ \hline
\id{L-04} & El componente empírico está registrado y diseñado, pero no ejecutado. & El protocolo consta en OSF con fecha anterior a cualquier recolección. Los resultados corresponden formalmente a la Entrega 4. \\ \hline
\id{L-05} & Los componentes de IA no cuentan con conjunto de datos etiquetado del propio cultivo. & Los valores de \id{RNF-01} y \id{RNF-02} se fijaron a partir de la literatura del dominio \cite{suharjito2021,goh2025,rosbi2024} y se declaran como objetivo, no como resultado medido. \\ \hline
\id{L-06} & Las sesiones de validación \emph{walkthrough} realizadas son inferiores a las tres exigidas. & Se declara el número efectivo con su acta correspondiente; las restantes se ejecutarán antes de la Entrega 4. \\ \hline
\id{L-07} & \emph{Resuelta.} La obtención de la hoja de liquidación semanal (\id{EV-13}) incorporó un segundo tipo documental, distinto del plan de labores. & Se dispone de dos tipos documentales de la organización. La obtención de un tercer tipo (hoja de conteo del polinizador o boleta de pesaje) queda como trabajo pendiente. \\ \hline
\id{L-10} & El prototipo alcanza una cobertura del 47{,}4\% de los requisitos \emph{Must}, por debajo del 60\% exigido por la §3.6 de la guía. & Se declara el valor real verificado por inspección del código, sin computar las funcionalidades simuladas como implementadas. Plan de incorporación en la §7.5. \\ \hline
\id{L-11} & El prototipo carece de servicio de respaldo: la persistencia es local al navegador y las credenciales se almacenan sin derivación de clave. & Los requisitos \id{RNF-10}, \id{RNF-11}, \id{RNF-13} y \id{RNF-14} no son verificables sobre el prototipo actual. Se declara en la §6.4 como desviación \id{D-02}. \\ \hline
\id{L-08} & Seis requisitos no funcionales y una restricción carecen de trazabilidad a evidencia de campo. & Se declaran en la §5.4.2 como derivados de decisiones de arquitectura y de prácticas estándar, sin atribuirles una fuente de elicitación. \\ \hline
\id{L-09} & El conflicto de interés por consanguinidad entre el analista líder y dos participantes persiste en la segunda ronda. & Se mantienen las seis medidas de mitigación declaradas en el Apéndice A, incluida la triangulación obligatoria y la prohibición de sustentar cualquier requisito exclusivamente en esas fuentes. \\ \hline
\caption{Limitaciones declaradas de la Entrega 3 (2A)}
\label{tab:limitaciones}
\end{longtable}

\subsection{Amenazas a la validez}

Siguiendo la clasificación de Wohlin \emph{et al.} \cite{wohlin2012}, se identifican las siguientes
amenazas al proceso de ingeniería de requisitos de este proyecto.

\textbf{Validez de constructo.} El cuestionario mide \emph{utilidad percibida} de funciones descritas
verbalmente, no de funciones utilizadas. La correlación entre utilidad declarada y adopción efectiva
no está establecida. Mitigación: las valoraciones se contrastaron con las dificultades declaradas de
forma independiente en la misma aplicación, y con las necesidades expresadas en las entrevistas
directivas; la convergencia entre las tres fuentes refuerza la interpretación.

\textbf{Validez interna.} Las entrevistas fueron conducidas por integrantes del equipo, con dos
participantes vinculados por consanguinidad al analista líder. Existe riesgo de sesgo de
complacencia. Mitigación: uso estricto del banco de preguntas aprobado, presencia de un segundo
integrante como registrador, disponibilidad de las grabaciones para auditoría, y verificación de
requisitos por integrantes sin participación en esas sesiones.

\textbf{Validez externa.} Los hallazgos provienen de una única unidad productiva de aproximadamente
cien hectáreas, con una variedad híbrida específica y en etapa de primera cosecha. No son
generalizables a plantaciones de mayor escala, de variedad \emph{guineensis} o en etapa madura. Los
umbrales de decisión se parametrizaron precisamente para acotar esta amenaza.

\textbf{Validez de conclusión.} El componente empírico no fue ejecutado; no se formula ninguna
conclusión estadística en este documento. Las cifras del cuestionario se presentan de forma
descriptiva, con su $n$ y su porcentaje de no respuesta explícitos.

\subsection{Trabajo pendiente para la Entrega 4 (2B)}

\begin{longtable}{|C{1.4cm}|L{7.0cm}|L{6.6cm}|}
\hline
\rowcolor{grisclaro}\textbf{Prioridad} & \textbf{Actividad} & \textbf{Resultado esperado} \\
\hline
\endhead
Alta & Ejecutar el componente empírico registrado en OSF con al menos tres personas evaluadoras independientes & Resultados con prueba de hipótesis, tamaño del efecto y acuerdo inter-evaluador \\ \hline
Alta & Completar las sesiones de validación \emph{walkthrough} con actas firmadas & Cobertura de las seis sesiones acumuladas del objetivo 2B \\ \hline
Alta & Construir el conjunto de datos etiquetado del cultivo & Habilita la medición real de \id{RNF-01} y \id{RNF-02} \\ \hline
Media & Ampliar la ronda de entrevistas hasta el objetivo acumulado & Curva de saturación con evidencia de estabilización \\ \hline
Alta & Elevar la cobertura del MVP por encima del 60\% incorporando \id{RF-10}, \id{RF-26} y \id{RF-35} & Cumplimiento del mínimo de la §3.6; los tres son de implementación acotada \\ \hline
Alta & Corregir las desviaciones \id{D-01} a \id{D-03} del prototipo & Coherencia entre el prototipo y el dominio real documentado \\ \hline
Media & Completar la implementación de los componentes de IA y del conteo georreferenciado & Sustitución de las funcionalidades simuladas por procesamiento real \\ \hline
Media & Obtener un tercer tipo documental de la organización & Completa la triangulación documental (\id{L-07}) \\ \hline
Media & Depositar el conjunto de datos anonimizado en Zenodo con licencia CC BY 4.0 & DOI asignado conforme a los principios FAIR \cite{wilkinson2016} \\ \hline
Baja & Diseñar la consolidación multi-finca & Prepara la extensión del alcance declarada como \emph{Won't} en esta versión \\ \hline
\caption{Trabajo pendiente para la Entrega 4 (2B)}
\end{longtable}

\subsection{Consideración final}

El resultado más relevante de esta entrega no es el incremento cuantitativo de requisitos, sino la
incorporación de una fuente externa a la organización cliente. La planta extractora ---que compra la
producción, califica su calidad y aplica penalizaciones económicas--- define los criterios que
determinan el valor del producto, y no había sido consultada en las entregas anteriores. Su
incorporación reveló que varios requisitos previamente especificados estaban formulados desde la
perspectiva exclusiva del productor, sin considerar el criterio de quien recibe el fruto.

Este hallazgo confirma la observación de Ferrari \emph{et al.} \cite{ferrari2016} sobre el
conocimiento tácito en las entrevistas de elicitación: el administrador y el asesor técnico conocían
el criterio de desprendimiento, pero no lo enunciaron porque para ellos era evidente. Fue necesario
entrevistar a quien aplica el criterio para que este se hiciera explícito y verificable.
